% subsection content: 隣接3項間漸化式型 \, $\displaystyle a_{n+2}+pa_{n+1}+qa_n=0$

\subsubsection{二次特性方程式型 \, $\displaystyle a_{n+2}+pa_{n+1}+qa_n=0$}
節は変われど,特性方程式を使って考える.
この小節でも0除算が発生しうる状況があるが,そのような場合はそもそもそのような漸化式が成立しないか,別の解法で解くものとなっているため数学的な問題はない.
結局,大方の場合,等比数列に変形できればいいのだから,これを等比数列に変形できないかといくつかの試行錯誤を重ねるのが王道の解き方である.
結果から言ってしまえば,特性方程式型のところで少し触れた式を生み出すことで等比数列を作ることができる.
\[
  a_{n+2}+pa_{n+1}+qa_n=0
\]
について,
\[
  a_{n+2}+\alpha a_{n+1}=\beta(a_{n+1}+\alpha a_{n})
\]
となるような$\alpha,\beta$の組を探せば良い.
変形すると
\begin{align*}
   & a_{n+2}+\alpha a_{n+1}=\beta(a_{n+1}+\alpha a_{n})               \\
   & \therefore \quad a_{n+2}+(\alpha-\beta)a_{n+1}-\alpha\beta a_n=0
\end{align*}
である.ここで$\beta'=-\beta$と置き,不動点を考えているから$a_n$に関する項を$x$と置けば,
\[
  x^2+(\alpha+\beta')x+\alpha\beta'=0
\]
となり,まさに二次方程式の解を求めることに一致する.
従って形式的には
\[
  a_{n+2}+pa_{n+1}+qa_n=0
\]
については,
\[
  x^2+px+q=0
\]
について解くことで既知の形に帰着する.

\,$x^2+px+q=0$の2解を$\alpha,\beta$とすると,元の漸化式は
\[
  \begin{cases}
    a_{n+1}-\alpha a_n=\beta(a_n-\alpha a_{n-1}) \\
    a_{n+1}-\beta a_n=\alpha(a_n-\beta a_{n-1})
  \end{cases}
\]
という2つの等比数列を生む(添字は1つずれているが本質的には先程と同じ議論である).
すなわち,\,$b_n=a_{n+1}-\alpha a_n$は公比$\beta$の等比数列,\,$c_n=a_{n+1}-\beta a_n$は公比$\alpha$の等比数列になる.
$\alpha\ne\beta$であれば,この2つの式から$a_{n+1}$を消去して$a_n$について解けばよい.
また,\,$\alpha=\beta$(重解)の場合は等比数列が1つしか作れないため,\,$b_n=a_{n+1}-\alpha a_n$が公比$\alpha$の等比数列であることを使い,
両辺を$\alpha^{n+1}$で割って階差型に帰着させることになる.

\noindent \textbf{【例題】}
\begin{enumerate}[label=(\arabic*),parsep=3pt]
  \item $\displaystyle a_1=1,\,a_2=1,\, a_{n+2}=5a_{n+1}-6a_n$
  \item $\displaystyle a_1=1,\,a_2=4,\, a_{n+2}=4a_{n+1}-4a_n$
  \item $\displaystyle a_1=2,\,a_2=1,\, a_{n+2}=a_{n+1}+6a_n$
\end{enumerate}

\noindent\textbf{【方針】}
\begin{enumerate}[label=(\arabic*),parsep=3pt]
  \item 特性方程式$x^2-5x+6=0$の解は異なる2つの実数解です.素直に2本の等比数列を作りましょう.
  \item 特性方程式$x^2-4x+4=0$を解くと重解になります.等比数列は1本しか作れないことに注意しましょう.
  \item 特性方程式の解の1つが負になるパターンです.符号に注意して計算しましょう.
\end{enumerate}

\noindent\textbf{【解答】}
\begin{enumerate}[label=(\arabic*),parsep=3pt]
  \item $\displaystyle a_n=2^n-3^{n-1}$
  \item $\displaystyle a_n=n\cdot2^{n-1}$
  \item $\displaystyle a_n=(-2)^{n-1}+3^{n-1}$
\end{enumerate}

\noindent\textbf{【解法】}
\begin{enumerate}[label=(\arabic*),parsep=3pt]
  \item $\displaystyle a_1=1,\,a_2=1,\, a_{n+2}=5a_{n+1}-6a_n$\\
        特性方程式$x^2-5x+6=0$を解くと,\\
        $(x-2)(x-3)=0$より$\alpha=2,\beta=3$である.
        \begin{align*}
           & a_{n+2}-2a_{n+1}=3(a_{n+1}-2a_n) \\
           & a_{n+2}-3a_{n+1}=2(a_{n+1}-3a_n)
        \end{align*}
        なので,\,$b_n=a_{n+1}-2a_n,\ c_n=a_{n+1}-3a_n$と置くと,
        \begin{align*}
           & b_n=b_1\cdot3^{n-1},\quad b_1=a_2-2a_1=-1    \\
           & c_n=c_1\cdot2^{n-1},\quad c_1=a_2-3a_1=-2    \\
           & \therefore \quad b_n=-3^{n-1},\quad c_n=-2^n
        \end{align*}
        $b_n-c_n=(a_{n+1}-2a_n)-(a_{n+1}-3a_n)=a_n$であるから,
        \[
          a_n=b_n-c_n=-3^{n-1}-(-2^n)=2^n-3^{n-1}
        \]
  \item $\displaystyle a_1=1,\,a_2=4,\, a_{n+2}=4a_{n+1}-4a_n$\\
        特性方程式
        \[
          x^2-4x+4=0
        \]
        を解くと,
        \[
          (x-2)^2=0
        \]
        より$\alpha=2$である.
        \[
          b_n=a_{n+1}-2a_n
        \]
        と置くと,
        \begin{align*}
          b_{n+1}&=a_{n+2}-2a_{n+1}\\
          &=(4a_{n+1}-4a_n)-2a_{n+1}\\
          &=2(a_{n+1}-2a_n)\\
          &=2b_n
        \end{align*}
        また,
        \begin{align*}
          b_1=a_2-2a_1=2
        \end{align*}
        から,
        \[
          b_n=2\cdot2^{n-1}=2^n
        \]
        したがって
        \[
          a_{n+1}-2a_n=2^n
        \]
        である.両辺を$2^{n+1}$で割ると,
        \[
          \frac{a_{n+1}}{2^{n+1}}-\frac{a_n}{2^n}=\frac12
        \]
        ここで
        \[
          c_n=\dfrac{a_n}{2^n}
        \]
        と置くと,
        \begin{align*}
          &c_{n+1}=c_n+\frac12,\quad c_1=\frac{a_1}{2}=\frac12\\
          &\therefore \quad c_n=\frac12+(n-1)\cdot\frac12=\frac n2\\
          &\therefore \quad a_n=2^n\cdot c_n=n\cdot2^{n-1}
        \end{align*}
  \item $\displaystyle a_1=2,\,a_2=1,\, a_{n+2}=a_{n+1}+6a_n$\\
        特性方程式
        \[
          x^2-x-6=0
        \]
        を解くと,
        \[
          (x-3)(x+2)=0
        \]
        より$\alpha=3,\beta=-2$である.
        \begin{align*}
          &a_{n+2}-3a_{n+1}=-2(a_{n+1}-3a_n)\\
          &a_{n+2}+2a_{n+1}=3(a_{n+1}+2a_n)
        \end{align*}
        なので,
        \[
          b_n=a_{n+1}-3a_n,\ c_n=a_{n+1}+2a_n
        \]
        と置くと,
        \begin{align*}
          &b_n=b_1\cdot(-2)^{n-1},\quad b_1=a_2-3a_1=-5\\
          &c_n=c_1\cdot3^{n-1},\quad c_1=a_2+2a_1=5\\
          &\therefore \quad b_n=-5(-2)^{n-1},\quad c_n=5\cdot3^{n-1}
        \end{align*}
        \[
          c_n-b_n=(a_{n+1}+2a_n)-(a_{n+1}-3a_n)=5a_n
        \]
        であるから,
        \begin{align*}
          5a_n&=5\cdot3^{n-1}-\{-5(-2)^{n-1}\}\\
          &=5\cdot3^{n-1}+5(-2)^{n-1}\\
          \therefore\quad a_n&=(-2)^{n-1}+3^{n-1}
        \end{align*}
\end{enumerate}

\subsubsection{等差二次特性方程式型 \, $\displaystyle a_{n+2}+pa_{n+1}+qa_n=r$}
今度は右辺が0ではなく$r$となった.
つまり$r$だけズレたのである.
ズレたらもとに戻す,つまり不動点を探すのである.
\[
  b_n=a_n-\alpha
\]
となる$\alpha$を求めると,
\begin{align*}
   & (a_{n+2}-\alpha)+p(a_{n+1}-\alpha)+q(a_n-\alpha)=0 \\
   & a_{n+2}+pa_{n+1}+qa_n=\alpha+p\alpha+q\alpha       \\
   & \therefore \quad r=\alpha +p\alpha +q\alpha        \\
   & \Leftrightarrow \quad \alpha=\frac{r}{p+q+1}
\end{align*}
あとは前節で述べた隣接3項間漸化式の問題である.

\noindent \textbf{【例題】}
\begin{enumerate}[label=(\arabic*),parsep=3pt]
  \item $\displaystyle a_1=2,\,a_2=2,\, a_{n+2}=5a_{n+1}-6a_n+2$
  \item $\displaystyle a_1=1,\,a_2=2,\, a_{n+2}=4a_{n+1}-4a_n+3$
\end{enumerate}

\noindent\textbf{【方針】}
\begin{enumerate}[label=(\arabic*),parsep=3pt]
  \item まず不動点$\alpha$を求め,\,$b_n=a_n-\alpha$と置いて前節の二次特性方程式型に帰着させましょう.
  \item こちらは重解のパターンです.同じ手順で$\alpha$を求めてから前節の重解の場合の処理を行いましょう.
\end{enumerate}

\noindent\textbf{【解答】}
\begin{enumerate}[label=(\arabic*),parsep=3pt]
  \item $\displaystyle a_n=1+2^n-3^{n-1}$
  \item $\displaystyle a_n=3+2^{n-2}(3n-7)$
\end{enumerate}

\noindent\textbf{【解法】}
\begin{enumerate}[label=(\arabic*),parsep=3pt]
  \item $\displaystyle a_1=2,\,a_2=2,\, a_{n+2}=5a_{n+1}-6a_n+2$
        \begin{align*}
          &a_{n+2}=5a_{n+1}-6a_n+2\\
          &a_{n+2}-5a_{n+1}+6a_n=2
        \end{align*}
        ここで
        \[
          (a_{n+2}-\alpha)-5(a_{n+1}-\alpha)+6(a_n-\alpha)=0
        \]
        を満たす$\alpha$を求めると,
        \[
          \alpha=1
        \]
        $b_n=a_n-1$と置くと,
        \begin{align*}
           & b_{n+2}=5b_{n+1}-6b_n \\
           & b_1=1,\, b_2=1
        \end{align*}
        （答案中略。前節の例題(1)と同じ漸化式に帰着する.全く同じ計算を$b_n$に対して行う)
        \[
          b_n=2^n-3^{n-1}
        \]
        を得るから,
        \[
          a_n=1+b_n=1+2^n-3^{n-1}
        \]
  \item $\displaystyle a_1=1,\,a_2=2,\, a_{n+2}=4a_{n+1}-4a_n+3$
        \begin{align*}
          &a_{n+2}=4a_{n+1}-4a_n+3\\
          &a_{n+2}-4a_{n+1}+4a_n=3
        \end{align*}
        ここで
        \[
          (a_{n+2}-\alpha)-4(a_{n+1}-\alpha)+4(a_n-\alpha)=0
        \]
        を満たす$\alpha$を求めると,
        \[
          \alpha=3
        \]
        なので,\,$b_n=a_n-3$と置くと,
        \begin{align*}
           & b_{n+2}=4b_{n+1}-4b_n \\
           & b_1=a_1-3=-2          \\
           & b_2=a_2-3=-1
        \end{align*}
        特性方程式$x^2-4x+4=0$は重解$\alpha'=2$を持つから,\,$c_n=b_{n+1}-2b_n$と置くと,
        \begin{align*}
           & c_{n+1}=2c_n                       \\
           & c_1=b_2-2b_1=-1-2(-2)=3            \\
           & \therefore \quad c_n=3\cdot2^{n-1}
        \end{align*}
        したがって$b_{n+1}-2b_n=3\cdot2^{n-1}$である.両辺を$2^{n+1}$で割ると,
        \begin{align*}
           & \frac{b_{n+1}}{2^{n+1}}-\frac{b_n}{2^n}=\frac34          \\
           & f_n=\dfrac{b_n}{2^n} \quad \text{と置くと,}                       \\
           & f_{n+1}=f_n+\frac34,\quad f_1=\frac{b_1}{2}=-1           \\
           & \therefore \quad f_n=-1+(n-1)\cdot\frac34=\frac{3n-7}{4} \\
           & \therefore \quad b_n=2^n\cdot f_n=2^{n-2}(3n-7)
        \end{align*}
        よって,
        \[
          a_n=3+b_n=3+2^{n-2}(3n-7)
        \]
\end{enumerate}
